% Source: https://artofproblemsolving.com/community/c3h2265598_ap_thread_amc_practice_thread
% Stable identity: aops-topic-2265598-q03
The midpoints of the sides of a right triangle cut it into four congruent triangles, each similar to it with ratio $\dfrac{1}{2}$, hence each with $\dfrac{1}{4}$ of its area. So each shaded triangle has $\dfrac{1}{4}$ of the area of the triangle it was cut from, and the next step works inside a triangle (the one at $G$) that also has $\dfrac{1}{4}$ of that area.

\begin{asy}
usepackage("amsmath");
size(5.5cm);
defaultpen(fontsize(10pt));
pair C=(0,0), A=(0,6), G=(6,0);
pair B=(0,3), D=(3,0), J=(3,3), K=(3,1.5), E=(4.5,0), I=(4.5,1.5), L=(4.5,0.75), F=(5.25,0), H=(5.25,0.75);
fill(B--C--D--cycle,lightgray); fill(K--D--E--cycle,lightgray); fill(L--E--F--cycle,lightgray);
draw(A--C--G--cycle,linewidth(0.9));
draw(B--D--J--cycle); draw(K--E--I--cycle); draw(L--F--H--cycle);
draw((0,0.35)--(0.35,0.35)--(0.35,0));
label("$A$",A,N); label("$C$",C,SW); label("$G$",G,SE); label("$B$",B,W); label("$D$",D,S); label("$J$",J,NE);
label("$K$",K,W); label("$E$",E,S); label("$I$",I,NE); label("$L$",L,W,fontsize(8pt)); label("$F$",F,S,fontsize(8pt));
label("$6$",(0,4.5),W); label("$6$",(1.5,0),S);
\end{asy}

The big triangle has area $\dfrac{1}{2}\cdot6\cdot6=18$, so the shaded areas are $\dfrac{18}{4},\dfrac{18}{16},\dfrac{18}{64},\dots$, and after $100$ steps the total is
\[\dfrac{18}{4}\left(1+\dfrac{1}{4}+\cdots+\dfrac{1}{4^{99}}\right)=\dfrac{18}{4}\cdot\dfrac{1-\left(\dfrac{1}{4}\right)^{100}}{1-\dfrac{1}{4}}=6\left(1-\dfrac{1}{4^{100}}\right).\]
This is less than $6$ by an extremely small amount, so the total is nearest to $\boxed{\textbf{(A)}\ 6}$.
